\chapter{Abstract}

\begin{tabular}{@{} >{\bfseries}p{0.22\textwidth} p{0.75\textwidth} @{}}
\MakeUppercase{Keywords} &
Inverter-Based Resources, Line Differential Protection,
Distance Protection, Microgrid Protection, Digital Twin,
Physics-Informed Neural Network, Fault Classification,
Wide-Area Protection Coordinator, IEC~61850 GOOSE Security,
Adaptive Relay Coordination.
\end{tabular}
\par

\begin{doublespacing}

The global power grid is undergoing a fundamental structural transformation:
synchronous machines---whose high-inertia fault signatures underpinned a
century of deterministic relay design---are being displaced at accelerating
pace by Inverter-Based Resources (IBRs). Grid-Following (GFL) and
Grid-Forming (GFM) converters are voltage-controlled current sources whose
fault currents are limited to 1.1--2.0~pu by fast sub-cycle current
limiters, and whose positive-, negative-, and zero-sequence contributions
are determined by software control laws rather than by physical machine
reactance. This paradigm shift systematically invalidates the foundational
assumptions of legacy overcurrent (ANSI~51/67), impedance distance (ANSI~21),
and differential (ANSI~87) protection algorithms.

This thesis proposes a hierarchical, self-adaptive protection and control
architecture---termed the \textbf{Self-Adaptive Multi-Bus Protection (SAMBP)}
framework---for IBR-dominated distribution and sub-transmission networks.
The framework spans six technical tiers, each developed from first principles
and validated through Monte Carlo simulation, digital twin emulation, and
hardware-in-the-loop testing.

Line differential protection (Chapter~4) develops a family of adaptive 87L
algorithms. The physics-constrained adaptive threshold modulates the operating
slope~$k_m$ as a function of the real-time IBR current ratio
$k_\mathrm{ibr} = I_\mathrm{IBR}/I_\mathrm{rated}$, eliminating the static
bias margin that previously caused missed trips under weak-infeed conditions.
Supplementary harmonic-discrimination (87LH), negative-sequence (87LN), and
transient-differential (TDCS-87L) elements extend sensitivity across the full
spectrum of unsymmetrical and high-impedance fault types. Chapter~5 derives the
apparent impedance trajectory seen by a distance relay under IBR reactive current
injection and proposes a quadrilateral characteristic with an IBR reach-correction
term. A Monte Carlo framework quantifies the reliability gain over classical mho
elements at $\Delta P_\mathrm{detect} = +7.6$ percentage points.

Chapter~6 characterises the protection-relevant distinction between GFL and
GFM inverter modes and proposes a ROCOF-supervised mode-switch islanding
detection scheme that achieves isolation within 120~ms while remaining
compliant with IEEE~1547-2018 ride-through requirements. An adaptive
overcurrent setting-group scheme uses IEC~61850~GOOSE mode flags to
reconfigure relay parameters within 50~ms of a grid-to-island transition.

Chapter~7 introduces a non-blocking real-time digital twin architecture
that runs an Extended Kalman Filter (EKF) to jointly estimate the IBR
current ratio, line impedance, and fault location parameter from current
measurements alone. The EKF achieves steady-state RMSE of 0.024~pu in
$k_\mathrm{ibr}$, 0.5\% in $Z_\mathrm{line}$, and 1.1\% in fault location
at a 50~$\mu$s relay sampling cycle.

Chapter~8 develops a Physics-Informed Neural Network (PINN) fault classifier
that embeds three physics constraints directly in the loss function.
The PINN achieves 96.1\% five-class accuracy on synthetic IBR fault data
while guaranteeing zero physics-constraint violations. A CUSUM change-point
detector triggers selective online re-training in response to fleet
composition drift, recovering full accuracy within 21 decision cycles.

Chapter~9 integrates all components into a Wide-Area Protection Coordinator
(WAPC) that localises faults via PMU voltage-dip triangulation (94.5\%
accuracy), optimises relay settings through an IEC standard-inverse
CTI-graded engine, secures GOOSE channels with HMAC-SHA256 (0.36~ms overhead,
54.9\% risk reduction), and validates the complete system through a
100-scenario hardware-in-the-loop regression suite.

Collectively, this thesis contributes 18 original algorithms and architectures,
of which 6 are Tier-1 (novelty $= 5$, significance $= 5$ on independent
expert assessment). The proposed SAMBP framework provides the first complete,
closed-loop, self-adaptive protection architecture for mixed GFL/GFM power
systems, from local relay physics through centralised coordination and
cybersecurity.

\end{doublespacing}
